%-------------------------
% Resume in Latex
% Author : Harshibar
% Based off of: https://github.com/jakeryang/resume
% License : MIT
%------------------------

\documentclass[letterpaper,10pt]{article}

\usepackage{latexsym}
\usepackage[empty]{fullpage}
\usepackage{titlesec}
\usepackage{marvosym}
\usepackage[usenames,dvipsnames]{color}
\usepackage{verbatim}
\usepackage{enumitem}
\usepackage[hidelinks]{hyperref}
\usepackage{fancyhdr}
\usepackage[T2A]{fontenc}
\usepackage[utf8]{inputenc}
\usepackage[english,russian]{babel}
\usepackage{tabularx}
% only for pdflatex
% \input{glyphtounicode}

% fontawesome
\usepackage{fontawesome5}

% fixed width
\usepackage[scale=0.90,lf]{FiraMono}

% light-grey
\definecolor{light-grey}{gray}{0.83}
\definecolor{dark-grey}{gray}{0.3}
\definecolor{text-grey}{gray}{.08}

\DeclareRobustCommand{\ebseries}{\fontseries{eb}\selectfont}
\DeclareTextFontCommand{\texteb}{\ebseries}

% custom underilne
\usepackage{contour}
\usepackage[normalem]{ulem}
\renewcommand{\ULdepth}{1.8pt}
\contourlength{0.8pt}
\newcommand{\myuline}[1]{%
  \uline{\phantom{#1}}%
  \llap{\contour{white}{#1}}%
}
\newcommand{\latinmono}[1]{{\fontencoding{T1}\ttfamily #1}}


% custom font: helvetica-style with Cyrillic support
\usepackage[default,scale=0.92]{opensans}


\pagestyle{fancy}
\fancyhf{} % clear all header and footer fields
\fancyfoot{}
\setlength{\footskip}{4.1pt}
\renewcommand{\headrulewidth}{0pt}
\renewcommand{\footrulewidth}{0pt}

% Adjust margins
\addtolength{\oddsidemargin}{-0.5in}
\addtolength{\evensidemargin}{0in}
\addtolength{\textwidth}{1in}
\addtolength{\topmargin}{-.5in}
\addtolength{\textheight}{1.0in}

\urlstyle{same}

\raggedbottom
\raggedright
\setlength{\tabcolsep}{0in}

% Sections formatting - serif
% \titleformat{\section}{
%   \vspace{2pt} \scshape \raggedright\large % header section
% }{}{0em}{}[\color{black} \titlerule \vspace{-5pt}]

% TODO EBSERIES
% sans serif sections
\titleformat {\section}{
    \bfseries \vspace{0pt} \raggedright \large % header section
}{}{0em}{}[\color{light-grey} {\titlerule[2pt]} \vspace{-6pt}]

% only for pdflatex
% Ensure that generate pdf is machine readable/ATS parsable
% \pdfgentounicode=1

%-------------------------
% Custom commands
\newcommand{\resumeItem}[1]{
  \item\small{
    {#1 \vspace{-2pt}}
  }
}

\newcommand{\resumeSubheading}[4]{
  \vspace{-1pt}\item
    \begin{tabular*}{\textwidth}[t]{l@{\extracolsep{\fill}}r}
      \textbf{#1} & {\color{dark-grey}\small #2}\vspace{1pt}\\ % top row of resume entry
      \textit{#3} & {\color{dark-grey} \small #4}\\ % second row of resume entry
    \end{tabular*}\vspace{-4pt}
}

\newcommand{\resumeSubSubheading}[2]{
    \item
    \begin{tabular*}{\textwidth}{l@{\extracolsep{\fill}}r}
      \textit{\small#1} & \textit{\small #2} \\
    \end{tabular*}\vspace{-7pt}
}

\newcommand{\resumeProjectHeading}[2]{
    \item
    \begin{tabular*}{\textwidth}{l@{\extracolsep{\fill}}r}
      #1 & {\color{dark-grey}\small #2} \\
    \end{tabular*}\vspace{-4pt}
}

\newcommand{\resumeSubItem}[1]{\resumeItem{#1}\vspace{-4pt}}

\renewcommand\labelitemii{$\vcenter{\hbox{\tiny$\bullet$}}$}

% CHANGED default leftmargin  0.15 in
\newcommand{\resumeSubHeadingListStart}{\begin{itemize}[leftmargin=0in, label={}, itemsep=0pt, topsep=0pt]}
\newcommand{\resumeSubHeadingListEnd}{\end{itemize}}
\newcommand{\resumeItemListStart}{\begin{itemize}[leftmargin=0.18in, itemsep=0pt, topsep=1pt, parsep=0pt, partopsep=0pt]}
\newcommand{\resumeItemListEnd}{\end{itemize}\vspace{-2pt}}

\color{text-grey}

%-------------------------------------------
%%%%%%  RESUME STARTS HERE  %%%%%%%%%%%%%%%%%%%%%%%%%%%%


\begin{document}

%----------HEADING----------
\begin{center}
    \textbf{\Large Очкин Никита Валерьевич} \\ \vspace{5pt}
    \small \faPhone* \latinmono{+7 977 194-80-50} \hspace{1pt} $|$
    \hspace{1pt} \faEnvelope \hspace{2pt} \latinmono{nkt.ochkin@gmail.com} \hspace{1pt} $|$
    \hspace{1pt} \faGithub \hspace{2pt}
    \href{https://github.com/namenick91}{\myuline{\latinmono{github.com/namenick91}}} \hspace{1pt} $|$
    \hspace{1pt} \faMapMarker* \hspace{2pt} Москва
    \\ \vspace{-3pt}
\end{center}

%-----------EXPERIENCE-----------
\section{ОПЫТ РАБОТЫ}
  \resumeSubHeadingListStart

    \resumeSubheading
      {ИСП РАН}{янв. 2026 -- настоящее время}
      {Лаборант}{Москва}
      \resumeItemListStart
        \resumeItem{Оценка VLM OCR-модели по качеству и скорости; подбираю модель для рабочих процессов.}
        \resumeItem{Разработка FastAPI-сервиса поверх PaddleOCR-VL: инференс через vLLM, Docker-стек из custom-api, paddleocr-vl-api и paddleocr-vlm-server.}
        \resumeItem{Бенчмарк сервиса; сейчас исследую VLM для извлечения данных из графиков в табличный вид.}
      \resumeItemListEnd

    \resumeSubheading
      {МГТУ им. Н.Э. Баумана}{окт. 2024 -- настоящее время}
      {Техник, кафедра ФН-11, прикладная математика}{Москва}
      \resumeItemListStart
        \resumeItem{Программные и ML-задачи кафедры.}
        \resumeItem{Разрабатка RAG-системы для образовательной платформы NOMOTEX: ETL, эмбеддинги, retrieval и генерация на локальной LLM.}
      \resumeItemListEnd

  \resumeSubHeadingListEnd


%-----------RESEARCH-----------
\section{НАУЧНАЯ РАБОТА}
    \resumeSubHeadingListStart
      \resumeProjectHeading
          {\textbf{Convformer}: прогнозирование временных рядов}{2026}
          \resumeItemListStart
            \resumeItem{Готовлю препринт для публикации в журнале Q2-Q3; препринт: \href{https://github.com/namenick91/Convformer}{\myuline{GitHub}}.}
            \resumeItem{Расширил эмбеддинг Informer двухслойным сверточным блоком, добавил декомпозицию ряда из Autoformer и заменил ProbSparse на линейное внимание FAVOR+ (Performer).}
            \resumeItem{На длинных горизонтах получил меньший MSE на стандартных бенчмарках.}
          \resumeItemListEnd

    \resumeSubHeadingListEnd


%-----------HACKATHONS-----------
\section{ХАКАТОНЫ}
    \resumeSubHeadingListStart
      \resumeProjectHeading
          {\textbf{МТС True Tech Hack 2026}: GPTHub}{2026}
          \resumeItemListStart
            \resumeItem{Был капитаном команды из 5 человек.}
            \resumeItem{Собрали платформу на базе OpenWebUI: маршрутизация запросов, RAG, долгосрочная память на PostgreSQL + PGVector и Deep Research через gpt-researcher.}
            \resumeItem{Отвечал за архитектуру RAG, подбор моделей и пресетов, двухуровневую память, интеграцию Deep Research и координацию команды.}
          \resumeItemListEnd
    \resumeSubHeadingListEnd


%-----------PROJECTS-----------
\section{PET-ПРОЕКТЫ}
    \resumeSubHeadingListStart
      \resumeProjectHeading
          {\textbf{Customer Churn} (CatBoost)}{\href{https://github.com/namenick91/ML-Workshelf\#customer-churn-prediction}{\myuline{GitHub}}}
          \resumeItemListStart
            \resumeItem{Построил модель оттока клиентов: очистка и типизация данных, фичи, k-fold и подбор гиперпараметров в Optuna.}
          \resumeItemListEnd

      \resumeProjectHeading
          {\textbf{Multimodal RAG} для поиска фильмов}{\href{https://github.com/namenick91/rag-multimodal-movies}{\myuline{GitHub}}}
          \resumeItemListStart
            \resumeItem{Собрал FastAPI + Qdrant + MinIO + Streamlit; использовал E5/CLIP-эмбеддинги и локальную Qwen для суммаризации.}
          \resumeItemListEnd

      \resumeProjectHeading
          {\textbf{Simpsons Character Classification} (EfficientNetV2-S)}{\href{https://github.com/namenick91/ML-Workshelf\#simpsons-character-classification}{\myuline{GitHub}}}
          \resumeItemListStart
            \resumeItem{Обучил multi-class классификатор изображений: двухфазный fine-tuning, RandAugment, CutMix, MixUp, AdamW + Cosine и early stopping.}
          \resumeItemListEnd

      \resumeProjectHeading
          {\textbf{Semantic Segmentation} (U-Net / SegNet)}{\href{https://github.com/namenick91/ML-Workshelf\#semantic-segmentation}{\myuline{GitHub}}}
          \resumeItemListStart
            \resumeItem{Сравнил модели для бинарной сегментации медицинских изображений; считал BCE, Dice, Focal и IoU, останавливал обучение по валидации.}
          \resumeItemListEnd
    \resumeSubHeadingListEnd



%-----------EDUCATION-----------
\section{ОБРАЗОВАНИЕ}
  \resumeSubHeadingListStart
    \resumeSubheading
      {МГТУ им. Н.Э. Баумана}{2022 -- 2026}
      {Бакалавриат, математика и компьютерные науки}{Москва}
      	\resumeItemListStart
    	\resumeItem{Факультет фундаментальных наук; профиль: суперкомпьютерное моделирование и искусственный интеллект в инженерных задачах.}
        \resumeItemListEnd
  \resumeSubHeadingListEnd


%
%-----------PROGRAMMING SKILLS-----------
\section{НАВЫКИ}
 \begin{itemize}[leftmargin=0in, label={}]
    \small{\item{
     \textbf{Языки программирования}{: Python, SQL}\vspace{2pt} \\
     \textbf{ML / DL}{: PyTorch, scikit-learn, CatBoost, Hugging Face Transformers; классическое ML, deep learning, прогнозирование временных рядов}\vspace{2pt} \\
     \textbf{Инфраструктура}{: bash, Linux, Docker, FastAPI, vLLM, PostgreSQL + PGVector, Qdrant, MinIO}\vspace{2pt} \\
     \textbf{Языки общения}{: русский -- родной; английский -- свободное владение}
    }}
 \end{itemize}


%-------------------------------------------
\end{document}
